\chapter{Conclusion}
\label{ch:conclusion}

% Répondez à la question posée dans l'introduction. Séparez résultats établis,
% limites et perspectives ; n'introduisez pas ici un nouveau théorème ou une
% nouvelle expérience.

Nous avons défini l'accessibilité comme la clôture réflexive transitive de
l'\cref{eq:transitive-closure}, démontré sa transitivité au
\cref{thm:transitivity} et calculé cette relation avec
l'\cref{alg:warshall}. Le graphe de la \cref{fig:example-graph} et la matrice du
\cref{tab:reachability} ont fourni une vérification transparente
des définitions et de l'algorithme.

L'exemple laisse volontairement ouvertes de nombreuses questions : reconstruction
des chemins, algorithmes pour graphes creux ou mises à jour incrémentales. Son
objectif principal reste structurel : montrer comment relier question de
recherche, définitions, résultats, preuves, limites et références correctement
citées, sans adopter une mise en page de type éditeur.
